%auto-ignore
\documentclass[main.tex]{subfiles}
\graphicspath{{\subfix{figs/}}}

\begin{document}

\title{Boundary Attention: Learning curves, corners, junctions and grouping}

\author{Mia Gaia Polansky$^{1}$\thanks{Much of this work was done while the author was a student researcher at Google.} \quad Charles Herrmann$^2$
\quad  Junhwa Hur$^2$ \\
\quad  Deqing Sun$^2$
\quad  Dor Verbin$^2$ \quad  Todd Zickler$^2$}
\authorrunning{M. Polansky et al.}
\institute{$^1$Harvard University \hspace{1cm}
\quad 
$^2$Google}

\maketitle
% \maketeaser

%%%%%%%%% ABSTRACT
\begin{abstract}
We present a lightweight network that infers grouping and boundaries, including curves, corners and junctions. It operates in a bottom-up fashion, analogous to classical methods for sub-pixel edge localization and edge-linking, but with a higher-dimensional representation of local boundary structure, and notions of local scale and spatial consistency that are learned instead of designed. Our network uses a mechanism that we call boundary attention: a geometry-aware local attention operation that, when applied densely and repeatedly, progressively refines a pixel-resolution field of variables that specify the boundary structure in every overlapping patch within an image. Unlike many edge detectors that produce rasterized binary edge maps, our model provides a rich, unrasterized representation of the geometric structure in every local region.  We find that its intentional geometric bias allows it to be trained on simple synthetic shapes and then generalize to extracting boundaries from noisy low-light photographs.
\end{abstract}

%%%%%%%%% BODY TEXT
\section{Introduction}
\label{sec:intro}

Converting a precise contour that is defined mathematically in continuous 2D space to a discrete pixel representation is a common task in computer graphics, and there are established tools for rasterization \cite{foley1982fundamentals,pineda1988parallel}, anti-aliasing \cite{heckbert1989fundamentals,catmull1974subdivision} and so on. However, the inverse problem in computer vision of robustly inferring precise, unrasterized contours from discrete images remains an open challenge, especially in the presence of noise. 

Earlier work by Canny~\cite{canny1986computational} and many others~\cite{concetta1988feature,parent1989trace,freeman1991design,freeman1992steerable,martin2004learning} explored the detection of unrasterized parametric edge models, and there are a variety of bottom-up algorithms that try to connect them by encouraging geometric consistency via edge-linking or message-passing. But since the dawn of deep learning, boundaries have almost exclusively been represented using discrete, rasterized maps; and spatial-consistency mechanisms that were previously based on explicit curve geometry have largely been replaced by black-box neural networks. 

In this paper, we take inspiration from early computer vision work and revisit the task of finding unrasterized boundaries via bottom-up geometric processing. But unlike early work, we leverage self-attention to learn these processes instead of hand-crafting them, thereby combining the benefits of geometric modeling with the efficiency and representational power of deep learning.

We focus on the low-level task of finding boundaries that separate regions of uniform color, as depicted by the toy problem in Figure~\ref{fig:overview}. This task becomes difficult at high noise levels, especially within local patches, and we address it by creating a model that can learn to exploit a wide range of low-level cues, such as curves being predominantly smooth with low-curvature; curves tending to meet at corner and junction points; contours tending to have consistent contrast polarity throughout their extent; and colors tending to vary smoothly at locations away from boundaries.

The core of our model is a mechanism we call boundary attention. It is a boundary-aware local attention operation that, when applied densely and repeatedly, progressively refines a field of variables that specify the local boundaries within dense (stride-1) patches of an image. The model's output is a dense field of unrasterized geometric primitives that, as depicted in the right of Figure~\ref{fig:overview}, can be used in a variety of ways, including to produce an unsigned distance function and binarized map for the image boundaries, a boundary-aware smoothing of the input image colors, and a map of spatial affinities between any query point and the pixels that surround it.

An important feature of our model is that its output patch primitives (see Figure~\ref{fig:junction-space}) have enough flexibility to represent a wide range of local boundary patterns and scales, including thin bars, edges, corners and junctions, all with various sizes, and all without rasterization and so with unlimited resolution. This gives our model enough capacity to learn to localize boundaries with high precision, and to regularize them in the presence of noise without erroneously rounding their corners or missing their junction points.

\begin{figure*}[t]
\centering
\includegraphics[width=\textwidth]{./figs/process_overview.png}
\caption{\label{fig:overview} Pipeline overview. The image unfolds into stride-1 patches, and boundary attention operates iteratively on their embeddings to produce for each patch: (\emph{i}) a parametric three-way partitioning, and (\emph{ii}) a parametric windowing function that defines its effective patch size. (Figure~\ref{fig:junction-space} shows parameterization details.) This output field implies a variety of global maps, shown in clockwise order: a boundary-aware smoothing of the input colors; an unsigned boundary-distance map; a boundary map; and a map of spatial affinities between any query point and its neighbors.}
\end{figure*}

We intentionally design our model with components that are local and invariant to discrete spatial shifts, enabling it to be trained on small-sized images and then deployed on much larger and differently-shaped ones. By tying the weights across multiple portions of the network, we increase its resilience to noise, which we validate through ablations. Our resulting model is very compact, comprising only 207k parameters, and it runs much faster than comparable optimization-based approaches. Further, our model can be trained to a useful state with very simple synthetic data, made up of random circles and triangles that are uniformly colored and then corrupted by noise. Despite the simplicity of this training data, we find that the model's learned internal activations exhibit intuitive behavior, and that the model generalizes to real world photographs, including those taken in low-light and have substantial shot noise.

Our main contributions can be summarized as follows:

\begin{enumerate}
    \item We introduce a bottom-up, feedforward network that decomposes an image into an underlying field of local geometric primitives that explicitly identify curves, corners, junctions and local grouping. 
    \item We do this by introducing a new parameterization for local geometric primitives, and a new self-attention mechanism that we call boundary attention.
    \item We identify the best architecture among a family of candidates, and we show that it generalizes from simple synthetic images to real low-light photographs.
\end{enumerate}

\begin{figure}[t]
\centering
\includegraphics[width=.7\columnwidth]{./figs/junction_overview.png}
\caption{\label{fig:junction-space} Parameterization details. \emph{Left}: Each patch $k$ is associated with an unrasterized three-way partitioning of its area (colored blue, orange and purple here). The partitioning parameters comprise a vertex $(u,v)$, orientation $\theta$, and angles $(\omega_1,\omega_2,\omega_3)$, defined up to scale. \emph{A}: A walk through junction space by linearly interpolating between junctions is spatially smooth, and can represent edges, bars, corners, Y-junctions, T-junctions and uniform regions. \emph{B}: Each junction is modulated through a learned windowing function. The windowing parameters $\mathbf{p}=(p_1,p_2,p_3)$ are convex weights over a dictionary of binary pillboxes.
}
\end{figure}



\section{Related Work}
\label{sec:related-work}

Early approaches to edge detection rely primarily on low-level or local information. Canny~\cite{canny1986computational} and others (\eg,~\cite{concetta1988feature,parent1989trace,freeman1991design,freeman1992steerable}) use  carefully chosen filters, and later ones like pB~\cite{martin2004learning} and gPb~\cite{4587420} combine such filters using a handful of trainable parameters. These methods provide local estimates of straight-edge position and orientation and they can subsequently be filtered and joined using geometry-based processes such as edge-linking~\cite{canny1986computational}. However, they often struggle near corners and junctions due to the difficulty of designing filters for these more complicated structures. Structured edges~\cite{dollar2014fast} and other early learning-based methods~\cite{dollar2006supervised, lim2013sketch} are able to detect more-complicated structures, but they represent their local outputs using discrete, rasterized boundary maps that are not directly compatible with continuous boundary parameterizations or geometry-based linking. 

Recently, the field of junctions (FoJ)~\cite{verbin2021field} introduced a way to detect more-complicated local structures while also representing them parametrically, without rasterization. By using its unrasterized representations in a customized non-convex optimization routine, the method can detect curves, corners and junctions with high precision and with unprecedented resilience to noise. Our model is strongly inspired by this work, and it improves upon it by introducing an enhanced family of parameterized geometric primitives (Figure~\ref{fig:junction-space}) and by replacing its hand-designed objective and associated non-convex optimization routine with a learned model that is fast, feed-forward and differentiable. We also tackle the challenge of representing images that vary in scale across different areas: whereas FoJ requires setting a single patch size for the entire image, our model learns to choose patch sizes adaptively in response to the input image.

We emphasize that our work is very different from a recent trend in computer vision of using deep networks to detect boundaries that are semantically meaningful (\eg,~\cite{su2021pixel, pu2022edter, zhou2023treasure}). These models aim to identify boundaries between semantic concepts, such as people and animals, while ignoring boundaries that are inside of these concepts, such as stripes on a shirt. In contrast to this trend, we follow earlier work by focusing entirely on identifying boundaries from the bottom up, using low-level cues alone. Since our model is trained using simple synthetic data, it does not exploit object familiarity or learn to ignore intra-object contrast. This approach has advantages: It is not specialized to any predetermined set of semantic concepts or tasks, and instead of producing a pixelized boundary map, it produces a field of unrasterized primitives that provide better precision as well as richer information about local grouping. We leave for future work the exploration of how to combine our bottom-up model with other cues that are top-down and semantic.

\section{Representation}
\label{sec:representation}

Our system is depicted in Figure~\ref{fig:overview}. It uses neighborhood cross-attention, a patch-wise variant of cross-attention, with $D$-dimensional, stride-1  embeddings. Critically, each $D$-dimensional embedding explicitly encodes a tuple of values that specifies the geometric structure and spatial extent of the unrasterized local boundaries within a patch. 

Our model relies on a  learned (linear) mapping from the embedding dimension $D$ to a hand-crafted, lower-dimensional space of unrasterized boundary patterns that we call \emph{junction space}. Junction space has the benefit of specifying per-patch boundary patterns without rasterization and thus with unlimited spatial precision. As depicted in Figure~\ref{fig:junction-space} and described in~\cite{verbin2021field}, it also has the benefit of including a large family of per-patch boundary patterns, including uniformity (\ie, absence of boundaries), edges, bars, corners, T-junctions and Y-junctions. Different from~\cite{verbin2021field}, we additionally  modulate the spatial extent of each patch by learning an associated windowing function, to be described in Section~\ref{sec:primitives}.

To enable communication between patch embeddings, each of which corresponds to a local patch, we leverage the idea that each image point is covered by multiple patches, and that overlapping patches must agree in their regions of overlap. This has the effect of tying neighboring patches together, analogous to cliques in a Markov random field. 
%Each patch's embedding should aptly represent its local geometry while agreeing with the structural claims of neighboring patches.

We refer to the core mechanism of our model as \emph{boundary attention}. We introduce it by first defining our hand-crafted parameterization of junction space and some associated operators. Then Section~\ref{sec:network} describes the architecture that we use to transform an image into a coherent junction-space representation.

\subsection{Boundary Primitives}
\label{sec:primitives}

We use parentheses $(x)$ for continuous signals defined on the 2D image plane $[0,W]\times[0,H]$ and square brackets $[n]$ for discrete signals defined on the pixel grid. We use $c[n]$ for the coordinates of the pixel with integer index $n$.

Denote the $C$-channel input image by $\{\mathbf{f}[n]\}$, where $\mathbf{f}[n]\in\mathbb{Q}^C$ is the vector image value at the discrete pixel grid index $n$. Our approach is to treat the image as a field of dense, stride-1 overlapping local patches. We use $\Omega_{k}(x)$ to denote the spatial support of the patch that is centered at the pixel whose integer index is $k$.

There are many ways to partition a local patch $\Omega_k(x)$, and one can define parametric families of partitions. For example the set of oriented lines provides a two-parameter family of partitions, with each member of the family separating the region into points that lie on one side of a line or the other. This family of partitions would be appropriate for describing edges. Here we define a larger family of partitions that encompasses a greater variety of local boundary structures.

As depicted in the right of Figure~\ref{fig:junction-space}, our partitions are parameterized by $\mathbf{g}\in\mathbb{R}^2\times\mathbb{S}^1\times\triangle^{2}$, where $\mathbb{S}^1$ is the unit circle and $\triangle^{2}$ is the standard $2$-simplex. We use the notation $\mathbf{g}=(\boldsymbol{u}, \theta, \boldsymbol{\omega})$, where $\boldsymbol{u}=(u,v)\in \mathbb{R}^2$ is the \emph{vertex}, $\theta\in\mathbb{S}^1$ is the \emph{orientation}, and   $\boldsymbol{\omega}=(\omega_1,\omega_2,\omega_3)$ are barycentric coordinates (defined up to scale) for the three relative angles, ordered clockwise starting from $\theta$. Our convention is to express the vertex coordinates relative to the center of region $\Omega_k(x)$, which is located at $c[k]$, and we note that the vertex is free to move outside of this region. We also note that up to two angles $\omega_j$ can be zero. This all makes it possible to smoothly represent a variety of partition types, including edges, bars, corners, 3-junctions and uniformity (\ie, trivial or singleton partitions), and do so with unlimited spatial resolution. 

Fixing a value for $\mathbf{g}$ induces three wedge support functions, denoted by
\begin{equation}
s_{kj}(x;\mathbf{g})\in\{0,1\},\ j=1, 2, 3.
\end{equation}
These evaluate to $1$ for points $x$ that are in $\Omega_k(x)$ and in the $j$th wedge defined by $\mathbf{g}$; and $0$ otherwise. It also induces an unsigned distance function, denoted by
\begin{equation}
d_{k}(x;\mathbf{g})\ge 0,
\end{equation}
which represents the Euclidean distance from point $x$ to the nearest point in the boundary set defined by $\mathbf{g}$. Figure~\ref{fig:junction-space} uses three colors to visualise the wedge supports $s_{kj}$ of a junction $\mathbf{g}$, and Figure~\ref{fig:outputs} shows the unsigned distance functions for $3\times 3$ grids of junction parameters. Analytic expressions for them are included in the supplement.

In order to enable the size of each region $\Omega_k$ to adapt to the local geometry and noise conditions, we equip each one with a parameterized local windowing function $w_k(x; \mathbf{p})\in[0,1]$, with parameters $\mathbf{p}\in{\cal P}=\triangle^{W-1}$ that are the coefficients of a convex combination of $W$ square pillbox functions. That is,
\begin{equation} \label{eq:windowfnconcrete}
    w_k(x;\mathbf{p}) = \sum_{i=1}^W p_i \mathbf{1}[\|x - c[k]\|_\infty \leq D_i],
\end{equation}
where $\|\cdot\|_{\infty}$ is the $\ell^{\infty}$-norm, and $\mathbf{1}[\cdot]$ is the indicator function that returns $1$ if the argument is true; and $0$ otherwise. In our experiments we use $W=3$ and diameters $\mathbf{D} = (3, 9, 17)$. \Cref{fig:outputs} shows some examples.


\begin{figure}[t]
\center
\includegraphics[width=\columnwidth]{./figs/output_visualization.pdf}
\caption{\label{fig:outputs} Example of our model's output, with examples from two different regions. \emph{Top row}: Some of each region's overlapping input patches, their corresponding outputs (visualized in the style of Figure~\ref{fig:junction-space}), and three types of per-patch attributes that the outputs imply: unsigned distance; boundaries; and gathered wedge colors. \emph{Bottom row}: Four types of global maps that are implied by accumulating values from the output field and rendered patches.}
\end{figure}

\subsection{Gather and Slice Operators}

Our network operates by refining the field $\{(\mathbf{g}^t[k], \mathbf{p}^t[k])\}$ over a fixed sequence of steps $t=1,...,T$. It uses two operators that we define here and depict in the right of Figure~\ref{fig:network} to facilitate intra-patch communication between pixels within a single patch, and inter-patch communication across different patches representing overlapping image regions. The first operator is a patch-wise \emph{gather} operator, in which each wedge of each patch computes the weighted average of the image values it contains (recall that $c[n]$ are the $n$th pixel's coordinates, and wedges are indexed by $j$):
\begin{equation}
    \mathbf{f}_{kj}=\frac{\sum_n \mathbf{f}[n]w_{k}(c[n];\mathbf{p}[k])s_{kj}(c[n];\mathbf{g}[k])}{\sum_n w_{k}(c[n];\mathbf{p}[k])s_{kj}(c[n];\mathbf{g}[k])}.\label{eq:gather}
\end{equation}

The second operation is a global pixel-wise \emph{slice} operation, where each pixel computes the means and variances, over all regions that contain it, of the  per-region distance maps $d_k(x;\mathbf{g}[k])$ and gathered wedge features $\mathbf{f}_{kj}$. The expressions for the means are:

\begin{align}
    \bar{d}[n] &=\frac{\sum_{k} w_{k}(c[n];\mathbf{p}[k])d_{k}(c[n];\mathbf{g}[k])}{\sum_{k} w_{k}(c[n];\mathbf{p}[k])}, \label{eq:slice-d}\\
    \bar{\mathbf{f}}[n] & =\frac{\sum_{k} w_{k}(c[n];\mathbf{p}[k])\sum_j \mathbf{f}_{kj}s_{kj}(c[n];\mathbf{g}[k])}{\sum_{k} w_{k}(c[n];\mathbf{p}[k])\sum_j s_{kj}(c[n];\mathbf{g}[k])}.\label{eq:globalf}
\end{align}
Here, the sums are over $\{k \mid \Omega_{k}\ni c[n]\}$ so that only patches that contain $c[n]$ contribute to the sum. Similar expressions for pixel-wise distance map variance $\nu_d[n]$
and feature variance $\nu_f[n]$, which is computed across patches containing $n$ and across their $K$ channels, are included in the supplement. Intuitively, slicing represents an accumulation of regional information into a pixel, as dictated by the partitions of all of the patches that contain the pixel.

\subsection{Visualizing Output}
\label{sec:output}

Our network's output is a field of tuples representing the junction and windowing parameters $\{(\mathbf{g}[k], \mathbf{p}[k]\}$ for all stride-1 patches of the input image. We  visualize them in Figure~\ref{fig:outputs} by rasterizing the continuous windowing functions $w_k(x;\mathbf{p})$ (second column) and binary wedge supports $s_{kj}(x;\mathbf{g})$, which are colored purple, orange, and blue (third column). To the left, we show the input image and the regions from which the patches were extracted.

Additionally, we can use the output junction parameters to rasterize unsigned distance patches $d_k(x;\mathbf{g})$ (fourth column), boundary patches $b_k(x;\mathbf{g})$ (fifth column), and wedge supports  $s_{kj}(x;\mathbf{g})$ that are re-colored with their respective wedge features $\mathbf{f}_{kj}$ (sixth column). Note that all of these are defined continuously, so can be rasterized at any resolution.

We expect the shapes of junction boundaries in overlapping regions to agree, so that the variances $\nu_d[n], \nu_f[n]$ are small at every pixel. Then, the fields of means $\{\bar{d}[n]\}, \{\bar{\mathbf{f}}[n]\}$ can be interpreted, respectively, as a global unsigned distance map for the image boundaries  (bottom row, fourth column) and a boundary-aware smoothing (bottom row, last column) of its input channel values. 
Figure~\ref{fig:outputs} shows an example (bottom row, fifth column), where we visualize the zero-set of the global unsigned distance map---we call this the global boundaries---by gathering boundary patches $b_k$ defined as:
\begin{equation} \label{eq:nonlinearboundary}
    b_k(x;\mathbf{g})=\left(1 + (d_k(x;\mathbf{g})/\eta)^{2}\right)^{-1},
\end{equation}
setting $\eta=0.3$.

For any query pixel $k$, we can also probe the containing wedge supports $\{s_{kj}(\cdot;\mathbf{g}[k])\}$ and windowing functions $\{w_k(\cdot, \mathbf{p}[k])\}$ to compute a spatial affinity map $a_{k}(x)$ that surrounds the query pixel. This represents the affinity between point $c[k]$ and a surrounding neighborhood with diameter that is twice that of $\Omega(x)$. It is also the boundary-aware spatial kernel that turns the neighborhood of input features $\{\mathbf{f}[\cdot]\}$ into the gathered value $\bar{\mathbf{f}}[n]$ via \begin{equation} \label{eq:attentionmap}
%a_{k}(x) =\sum_{k:\ \Omega_{k}\ni x} w_{k}(x;\mathbf{p}[k])\sum_j s_{kj}(x;\mathbf{g}[k]).
%a_{k}(x) = w_{k}(x;\mathbf{p}[k])\sum_j s_{kj}(x;\mathbf{g}[k]).
\bar{\mathbf{f}}[n] = \sum_{k} a_n(c[k]) \mathbf{f}[k].
\end{equation}
The expression for $a_n(x)$ follows from inserting  Equation~\ref{eq:gather} into~\ref{eq:globalf}. The spatial affinity maps for two probed points are shown in the leftmost image of the bottom row of Figure~\ref{fig:outputs}. Like slicing, querying the affinity map at a pixel is a form of regional accumulation from patches to a pixel.

% \clearpage
\begin{figure}[t]
\centering
\includegraphics[width=\columnwidth]{./figs/network.pdf}\caption{\label{fig:network} Model Architecture. All blocks are invariant to discrete spatial shifts, and only colored blocks are learned. Orange blocks operate at individual locations $n$, while blue ones operate on small spatial neighborhoods. Symbol $\oplus$ is concatenation, and gather and slice operators (Eqs.~\ref{eq:gather}--\ref{eq:globalf}) are depicted at right. The first iteration uses $\boldsymbol{\gamma}^0[n]=\boldsymbol{\gamma}_0[n]$, $\bar{\mathbf{f}}^0[n]=\mathbf{f}[n]$, and $\boldsymbol{\pi}^0[n]=\boldsymbol{\pi}_o$ with $\boldsymbol{\pi}_o$ learned across the training set. Boundary attention repeats $T=8$ times, with one set of weights for the first four iterations and a separate set of weights for the last four iterations, resulting in 207k trainable parameters total.}
\end{figure}

\section{Network Architecture}
\label{sec:network}

We seek a differentiable architecture that can effectively initialize and then refine the fields $\{(\mathbf{g}^t[k],\mathbf{p}^t[k])\}$ using a manageable number of iterations. This approach is motivated by the edge localization and edge-linking steps of early edge detectors like Canny's, but has a more sophisticated local boundary model, and mechanisms for spatial consistency that are learned instead of designed. It is also analogous to the original field of junctions algorithm~\cite{verbin2021field}, which uses coordinate descent for initialization and iterations of gradient descent for refinement. We want to replace both steps with something that is differentiable, faster, and able to scale to larger images.

After considering a variety of alternatives, we settled on a particular family of architectures based on dot-product self-attention; and within this family, we performed an extensive set of ablations to determine the best performing model. We describe our final model here, and we provide the ablation details in Section~\ref{sec:ablations}. Importantly, all of our model's components are invariant to discrete spatial shifts of the image, operating either on individual locations $k$, or on small neighborhoods of locations with spatially-shared parameters. This means that our model can be trained on small images and then deployed on much larger ones. Also, our model has only $207,000$ learnable parameters, making it orders of magnitude smaller than most deep semantic boundary detectors. As a point of reference, Diffusion Edge, a recent diffusion-based semantic boundary detection model uses on the order of 300 million parameters~\cite{ye2024diffusionedgediffusionprobabilisticmodel}.

Our model is depicted in Figure~\ref{fig:network}. It represents the field elements as higher-dimensional embeddings $\boldsymbol{\gamma}^t[k]\in\mathbb{R}^{D_\gamma}$ and $\boldsymbol{\pi}^t[k]\in\mathbb{R}^{D_\pi}$, which can be decoded at any iteration using the learned linear mappings $\boldsymbol{\gamma}\mapsto \mathbf{g}$ and $\boldsymbol{\pi}\mapsto \mathbf{p}$. Our final model uses $D_\gamma = 64$ and $D_\pi = 8$ which we show in our experiments provides enough capacity to learn a smooth latent representation of junction space. Using separate embeddings for the junction and windowing fields provides a disentangled representation of both.

Given an input image, the network first applies a ``neighborhood MLP-mixer'', which is a modified MLP-Mixer~\cite{tolstikhin2021mlpmixer} that replaces the global spatial operations with convolutions of kernel size $3$. The other change is that we map the input pixels to the hidden state size with a pixel-wise linear mapping rather than taking patches of the input. This block transforms the input image into a pixel-resolution map with ${D_\gamma}$ channels. We denote this by $\gamma_0[n]$ and refer to it as the initial ``hidden state''. This hidden state is then refined by a sequence of eight boundary attention iterations, which we describe next. (See our experiments for a visualization of the decoded hidden states as they evolve.)

The eight iterations of refinement are broken into two blocks, each with learned weights. In each iteration, we first add a linear mapping of the initial hidden state to the current hidden state, which acts as a skip connection. Next, we clone our hidden state, concatenating a dimension 8 learned windowing embedding to one of the copies and the input image plus the current estimate of the smoothed global features to the other. We treat the copies as the inputs to neighborhood cross-attention: each pixel in the first copy does two iterations of cross attention with a size 11 patch of the second copy. We add a learned $11 \times 11$ positional encoding to the patch, which allows our network to access relative positioning, even though global position cues are absent. We follow each self attention layer with a small MLP.

To transform our output or intermediary hidden state into junction space and render patch visualizations, we use a simple linear mapping. We separate the windowing embedding (the last 8 dimensions) from the junction embedding (the first 64 dimensions) and project each through a linear layer. We map the junction embedding to 7 numbers that represent $\mathbf{g}=(\boldsymbol{u}, \sin(\theta), \cos(\theta), \boldsymbol{\omega})$. These serve as the inputs to our gather and slice operators. 

To extract $\mathbf{p}$ from the windowing embedding, we linearly project the windowing embedding to a length 3 vector, which we use as coefficients in a weighted sum of three square pillbox functions with widths $3$, $9$, and $17$. This implementation of the windowing function ensures spatial overlap between neighboring patches by limiting the minimum patch size to $3$. 

\subsection{Training}
\label{sec:training}
Estimating the best junctions for noisy image patches is a non-convex optimization that is prone to getting stuck in local minima~\cite{verbin2021field}. We find that training our network in three stages with increasingly complex synthetic training data produces the best model. First we train the network on small $21 \times 21$ images containing a single junction corrupted by low amounts of Gaussian noise. Upon convergence, we retrain the network on larger $100 \times 100$ images containing a single circle and triangle corrupted by moderate Gaussian noise. Finally, we retrain the network on a high-noise synthetic dataset containing $125 \times 125$ images consisting of many overlapping triangles and circles. (See supplement for examples.)

We calculate our loss using the outputs of the two final iterations of our network, where the loss of the final output is weighted three times as much as the loss applied to the output before it. This encourages the network to allocate capacity to producing high quality outputs, while still providing supervision to intermediate junction estimates.

We train our model using a combination of four global losses applied to global (\ie sliced) fields, and two patch-wise losses applied to individual patches. The first two losses are supervision losses penalizing mismatches between our network's predictions and the ground truth feature and distance maps:
\begin{align}
    \mathcal{L}_f &= \sum_{n} \alpha[n] \|\bar{\mathbf{f}}[n] - \mathbf{f}_{\text{GT}}[n]\|^2, \label{eq:floss}\\
    \mathcal{L}_d &= \sum_{n} \alpha[n] \left(\bar{d}[n] - d_{\text{GT}}[n]\right)^2,
\end{align}
where $\mathbf{f}_{\text{GT}}$ and $d_{\text{GT}}$ are the ground truth features and distance maps, respectively, and $\alpha[n]$ is a pixel importance function defined as
\begin{equation} \label{eq:alpha}
    \alpha[n] = e^{-\beta \cdot (d_{\text{GT}}[n] + \delta)} + C,
\end{equation}
with $\beta$ and $C$ controlling how much weight to give pixels near boundaries. We set $\beta = 0.1$, $\delta = 1$, and $C=0.3$. (The supplement contains additional tests with a more involved importance mask.)
Using noiseless feature maps for supervision in Equation~\ref{eq:floss} has the effect of encouraging windowing functions to be as large as possible, because larger regions imply averaging over a greater number of noisy input pixel values.

On top of the two supervision losses we apply two consistency losses borrowed from~\cite{verbin2021field}, which minimize the per-pixel variances $\nu_f[n]$ and $\nu_b[n]$. Here, we weight them by $\alpha[n]$ from Equation~\ref{eq:alpha}. These consistency losses encourage the junction shapes $\mathbf{g}$ in overlapping regions to agree. Minimizing $\nu_f[n]$ also provides a second mechanism to encourage windowing functions to be large. Larger windows increase gather area, thereby reducing noise in wedge features $\mathbf{f}_{nj}$ that are sliced to compute variance $\nu_f[n]$ at each pixel $n$.

Finally, we use two patch-wise losses to encourage individual  feature and distance patches to agree with the supervisory ones:
\begin{align}
    \ell_f = \sum_k \chi[k] \sum_{n \in \Omega_k} \alpha[n] \|\bar{\mathbf{f}}[n] - \mathbf{f}_\text{GT}[n]\|^2, \\
    \ell_d = \sum_k \chi[k] \sum_{n\in\Omega_k} \alpha[n] (\bar{d}[n] - d_{\text{GT}}[n])^2,
\end{align}
where $\chi[k]$ is a patch importance function defined as:
\begin{equation}
    \chi[k] = \left(\sum_{n\in\Omega_k}(d_\text{GT}[n] + \delta')\right)^{-1},
\end{equation}
with $\delta' = 1$. These per-patch losses provide a more direct signal for adjusting model weights, compared to per-pixel losses which average over multiple patches. 

\begin{figure}[t]
\center
\includegraphics[width=.9\columnwidth]{./figs/synthetic_results.pdf}
\caption{\label{fig:syntheticval} \textit{Left:} ODS F-score for our method and multiple baselines at different noise levels computed on noisy synthetic data. The bottom inset show example patches at representative PSNR values. Our method outperforms all baselines at low noise and is better or competitive with other techniques at high noise. \textit{Right:} Comparing the F-score for different techniques with their runtime. Our method has the best average F-score while also being much faster than the second best method Field of Junctions.
}
\end{figure}

\begin{figure}[t]
\center
\includegraphics[width=.9\columnwidth]{./figs/eld_results.png}
\caption{\label{fig:ELD-performance} \emph{Left}: Qualitative comparison of our model versus other methods for a noisy crop from the ELD dataset~\cite{wei2020physicsbased}. See the supplement for more examples. Our method provides more detail than other techniques at low noise and is more robust to high levels of noise. \emph{Right}:
Repeatability of the estimated boundaries of crops from the ELD dataset over increasing noise. Over multiple levels of real sensor noise, our method is the most consistent in predicting edges.
}

\end{figure}

\section{Experiments}
\label{sec:experiments}

\vspace{-5pt}
\paragraph{Implementation details.}

In the final stage of training, we use noisy synthetic data of many randomly-colored overlapping triangles and circles. We render $240\times 320$ images containing $15$ to $20$ shapes each, but use $125\times 125$ crops for training. To those crops we add Gaussian or Perlin noise~\cite{perlin1985image}, and with probability $0.1$ we average over the color channels to produce grayscale inputs. Our dataset contains $10^5$ images, $90\%$ of which are used for training, and the rest for validation.

\vspace{-5pt}
\paragraph{Baselines.}

Since we focus on bottom-up edge and corner detection, we compare against other techniques with a similar focus and that, like us, do not train on large semantic datasets: Canny~\cite{canny1986computational}, gPb~\cite{4587420}, Structured Edges (SE)~\cite{dollar2014fast}, and Field of Junctions (FoJ)~\cite{verbin2021field}. 

\vspace{-5pt}
\paragraph{Quantitative results.}

For quantitative evaluation we cannot use semantic edge detection benchmarks like BSDS500~\cite{arbelaez2010contour}, because our model and the comparisons are designed to predict all edges, including those that are not semantically meaningful. We instead rely on synthetic data, where the ground truth edges can be determined with perfect precision, and inputs can be controllably noised. Our evaluation data comprises geometric objects and per-pixel additive Gaussian noise with a random variance.

Figure~\ref{fig:syntheticval} compares the performance and inference time of our method and baselines under different noise levels. The tuneable parameters for Field of Junctions were chosen to maximize its performance on noisy images with $17\times 17$ patches. Notably our method's adaptive windowing function gives it an edge compared to the Field of Junctions at low noise, enabling it to capture finer details, with only slightly worse performance under extreme noise conditions. Our method is also orders of magnitude faster than FoJ, as shown at right.

\paragraph{Qualitative results on real images.}
As shown in Figure~\ref{fig:ELD-performance}, despite being trained on synthetic data, our method can detect edges in real photographs with multiple levels of real sensor noise present in ELD~\cite{wei2020physicsbased}. Our method produces crisp and well-defined boundaries despite high levels of noise. The supplement includes additional examples that show that our method makes reasonable boundary estimates for other captured images.

\paragraph{Repeatability on real images.} We quantify each method's noise resilience on real low-light images by measuring the repeatability of its boundaries over a collection of scenes from the ELD dataset. For each scene, we run the model on the lowest-noise image and then measure, via ODS F-score, how much its predictions change with increasing noise level. Figure~\ref{fig:ELD-performance} shows the averages of these scores across the collection. Our model provides more consistent results for increasingly noisy images than other methods, in addition to capturing fine details that other methods miss.

\paragraph{Properties of learned embedding.} We find that our model learns a spatially smooth embedding $\boldsymbol{\gamma}\in\mathbb{R}^{D_\gamma}$ of junction space $\mathbf{g}\in {\cal G}$. In Figure~\ref{fig:interpolation} we generate equally-spaced samples $\boldsymbol{\gamma}_i$ by linearly interpolating from a particular $\boldsymbol{\gamma}_a$ to $\boldsymbol{0}$ and then to a particular $\boldsymbol{\gamma}_b$; and then to each sample we apply the learned embedding to compute and visualize the implied junction $\mathbf{g}_i$. We see that the embedding space is smooth, and interestingly, that it learns to associate its zero with nearly-equal angles and a vertex close to the patch center. For visual comparison, we show an analogous geometric interpolation in junction space ${\cal G}$ (see the supplement for expressions) from $\mathbf{g}_a$ to $\mathbf{g}_0\triangleq(\mathbf{0},0,\nicefrac{1}{3}\cdot\mathbf{1})$ and then to $\mathbf{g}_b$.

\begin{figure}[t]
\center
\includegraphics[width=.7\columnwidth]{./figs/interpolation.pdf}
\caption{\label{fig:interpolation} \emph{Top}: Linear interpolation in our network's learned embedding space $\mathbb{R}^{D_\gamma}$ from value $\boldsymbol{\gamma}_a$ to zero and then to $\boldsymbol{\gamma}_b$. \emph{Bottom}: A geometric interpolation in junction space $\mathbf{g}\in{\cal G}$ that passes through $\mathbf{g}_0=(\mathbf{0},0, \nicefrac{1}{3}\cdot\mathbf{1})$. The embedding has learned to be smooth and have an intuitive zero.}
\end{figure}

\begin{figure}[t]
\centering
\includegraphics[width=\textwidth]{./figs/evolution.png}
\caption{\label{fig:evolution} Evolution of boundaries during iterations, in reading order. Early  iterations  are  exploratory  and  unstructured, while later iterations feature consistent per-patch boundaries, resulting in clean average boundary maps.}
\end{figure}

\begin{table}[t]
    \centering
{\scriptsize
\centering
\setlength{\tabcolsep}{4pt}
\begin{tabular}{c c c c | c c c }
%\multicolumn{5}{c}{} & ODS F-score  \\  \hline
NM &BA-1 &BA-2 &Tied & \! Low Noise$\uparrow$  & \! Med-Noise$\uparrow$  & \! High Noise$\uparrow$  \\ \hline
\multicolumn{1}{l}{\cmark} &\xmark &\xmark & N/A & $0.355$   & $0.189$  & $0.179$  \\
\multicolumn{1}{l}{\cmark} &\cmark &\xmark &\xmark & $0.837$   & $0.582$  & $0.293$ \\
\multicolumn{1}{l}{\cmark} &\cmark &\cmark &\xmark & $\textbf{0.874}$   & $0.658$  & $0.327$ \\
\hline \multicolumn{1}{|l}{\cmark} &\cmark &\cmark &\cmark & $0.872$  & $\textbf{0.673}$  & \multicolumn{1}{c|} {$\textbf{0.348}$} \\ \hline
\end{tabular}
}
\caption{\label{table:ablation1} Ablations: impact of component combinations and weight-tying on F-score (higher is better). Our reported model is boxed.}
\end{table}

\paragraph{Evolution.} Figure~\ref{fig:evolution} shows an example of how the distance map $\bar{d}[n]$ evolves during refinement. Specifically, we visualize the result of slicing $b_k(x; \mathbf{g})$, to which we apply a non-linearity that amplifies less-prominent boundaries. We see that early iterations are exploratory and unstructured, and that later iterations agree.

\section{Ablations}\label{sec:ablations}

\Cref{table:ablation1} shows what happens when the initial per-patch Neighborhood MLP-mixer [NM] is used alone, versus combining it with one boundary attention block [BA-1] or two such blocks [BA-2]. Tying the attention weights across iterations provides a slight advantage at higher noise levels, with only a slight penalty in accuracy at lower noise. Our final model (boxed) provides the best overall performance, accepting slightly lower accuracy at low noise in exchange for better accuracy at higher noise levels.

\Cref{table:ablation2} compares many other variations of the boxed model from \cref{table:ablation1}, with asterisks indicating the model specifications used in the paper. \emph{Iterations} varies the number of attention-iterations within each block, and \emph{Neighborhood} varies the attention neighborhood size. \emph{MLP Input} varies the features concatenated to the hidden state prior to the pre-attention MLP and shows that replacing the gathered colors $\bar{\mathbf{f}}^t$ with a constant array or the input image values $\mathbf{f}$ performs worse. \emph{Windowing} shows that using fixed, square windowing functions $w_n(x)$ of any size performs equal (at low noise) or worse than inferring them adaptively. Somewhat intuitively, a small $9\times 9$ patch size performs well at low noise, but performance lags under noisy conditions where using a larger patch size increases the spatial extent of the communication across patches.

\begin{table}[t]
\scriptsize
\centering
\setlength{\tabcolsep}{6pt}
{
\centering
\begin{tabular}{l | l | c c c l}
\multicolumn{1}{l}{} & & \! Low Noise$\uparrow$   & \! Med-Noise$\uparrow$   & \! High Noise$\uparrow$   \\ \hline
\multirow{3}{*}{Iterations}
 &3 & $0.867$  & $0.647$  & $0.329$ \\
 &4* & $\textbf{0.872}$  & $\textbf{0.673}$  & {$\textbf{0.348}$} \\
 &5 & $0.871$  & $0.663$  & $0.337$ \\\hline 
\multirow{4}{*}{Neighborhood} 
% &$3\times 3$ & $0.835$   & $0.531$  & $0.240$  \\
 &$7\times 7$ & $0.869$   & $0.637$  & $0.313$  \\
 &$9\times 9$ & $0.871$   & $0.658$  & $0.325$  \\
 &$11 \times 11$* & $\textbf{0.872}$  & $\textbf{0.673}$  & $\textbf{0.348}$ \\
 &$13\times 13$ & $0.871$   & $0.655$  & $0.324$  \\ \hline
\multirow{3}{*}{MLP Input} & Constant & $0.870$   & $0.650$  & $0.317$  \\
 &Input features & $0.868$   & $0.657$  & $0.316$  \\
 &Avg.~features* & $\textbf{0.872}$  & $\textbf{0.673}$  & $\textbf{0.348}$ \\ \hline
\multirow{2}{*}{Windowing}
&Fixed & $\textbf{0.872}$  & $0.662$  & $0.328$ \\
 &Inferred* & $\textbf{0.872}$  & $\textbf{0.673}$  &$\textbf{0.348}$ \\
\end{tabular}
}
\caption{\label{table:ablation2} Ablations: variations of our reported model. We compare the number of iterations of boundary attention, the attention neighborhood size, the features concatenated to the hidden state as input to the MLP, and a fixed versus learned windowing function. Asterisks indicate our reported choices.}
\end{table}

\section{Conclusion}
We have introduced a differentiable model that uses boundary attention to explicitly reason about geometric primitives such as edges, corners, junctions, and regions of uniform appearance in images.
Our bottom-up, feedforward network can encode images of any resolution and aspect ratio into a field of geometric primitives that describe the local structure within every patch. This work represents a step to the goal of synergizing the benefits of low-level parametric modeling with the efficiency and representational power of deep learning.

%%%%%%%%% REFERENCES
\ifSubfilesClassLoaded{
\bibliographystyle{splncs04}
\bibliography{refs}
}{}


\end{document}